\documentclass[12pt,a4paper]{article}

% ── Packages ─────────────────────────────────────────────────────────────────
\usepackage[margin=2.5cm]{geometry}
\usepackage{amsmath, amssymb, amsthm}
\usepackage{booktabs}
\usepackage{longtable}
\usepackage{array}
\usepackage{xcolor}
\usepackage{listings}
\usepackage{hyperref}
\usepackage{graphicx}
\usepackage{float}
\usepackage{caption}
\usepackage{enumitem}
\usepackage{titlesec}
\usepackage{parskip}
\usepackage{lmodern}
\usepackage[T1]{fontenc}
\usepackage[utf8]{inputenc}

% ── Code listing style ────────────────────────────────────────────────────────
\definecolor{codeblue}{rgb}{0.13,0.29,0.53}
\definecolor{codegray}{rgb}{0.5,0.5,0.5}
\definecolor{codegreen}{rgb}{0.0,0.45,0.0}
\definecolor{backcolour}{rgb}{0.97,0.97,0.97}

\lstdefinestyle{python}{
    language=Python,
    backgroundcolor=\color{backcolour},
    commentstyle=\color{codegreen},
    keywordstyle=\color{codeblue}\bfseries,
    stringstyle=\color{orange},
    numberstyle=\tiny\color{codegray},
    basicstyle=\ttfamily\small,
    breaklines=true,
    captionpos=b,
    numbers=left,
    numbersep=6pt,
    showstringspaces=false,
    frame=single,
    framerule=0.4pt,
    rulecolor=\color{gray!40},
}

\hypersetup{
    colorlinks=true,
    linkcolor=codeblue,
    urlcolor=codeblue,
    citecolor=codeblue,
}

% ── Math helpers ──────────────────────────────────────────────────────────────
\newtheorem{definition}{Definition}
\newtheorem{remark}{Remark}

% ── Title ─────────────────────────────────────────────────────────────────────
\title{%
    \textbf{Pension Fund Efficiency Toolkit}\\[0.4em]
    \large Technical Explanation: Calculations, Data, Graphs, and Process Tests\\[0.2em]
    \normalsize Zimbabwe Self-Administered Pension Funds (2018--2022)
}
\author{Pension Efficiency Research Project}
\date{March 2026}

% ─────────────────────────────────────────────────────────────────────────────
\begin{document}
\maketitle
\tableofcontents
\newpage

% =============================================================================
\section{Overview: What We Are Doing and Why}
% =============================================================================

Zimbabwe's pension sector manages the retirement savings of hundreds of thousands of
workers regulated by the Insurance and Pensions Commission (IPEC).  The sector
operates under extreme macroeconomic stress: hyperinflation peaked above 500\% in
2020, multiple currency reforms were enacted between 2018 and 2022, and COVID-19
disrupted member contributions and investment markets simultaneously.

\textbf{The central question of this study is:} \emph{Among Zimbabwean
self-administered pension funds, which funds convert their resources into member
benefits most efficiently, and what structural or environmental factors explain
differences in efficiency?}

We answer this question using a two-stage empirical framework:

\begin{enumerate}[leftmargin=2cm]
    \item \textbf{Stage 1 -- Efficiency measurement (Data Envelopment Analysis).}
          DEA is a non-parametric frontier method that assigns each Decision-Making Unit
          (DMU, here: each pension fund) an efficiency score $\theta \in (0, 1]$ by
          comparing it against the best-practice frontier formed by the data itself.
          A score of $1.0$ means the fund is on the frontier (fully efficient);
          a score of $0.7$ means it could achieve the same outputs with only $70\%$ of
          its current inputs.

    \item \textbf{Stage 2 -- Determinant analysis (Random Forest regression).}
          Once bias-corrected efficiency scores are obtained, we regress them on six
          contextual variables --- inflation, exchange rate volatility, fund age,
          log fund size, expense ratio, and returns-to-scale classification --- to
          identify which factors drive efficiency differences across funds.
\end{enumerate}

The pipeline also performs:
\begin{itemize}
    \item \textbf{Simar--Wilson bootstrap} to correct the upward bias inherent in
          DEA estimators and to construct confidence intervals around each score.
    \item \textbf{Scale efficiency analysis} to classify each fund as operating under
          Increasing, Decreasing, or Constant Returns to Scale.
    \item \textbf{Peer analysis} to identify which frontier funds serve as benchmarks
          for each inefficient fund and to compute concrete input-reduction targets.
\end{itemize}

% =============================================================================
\section{Dataset: Column-by-Column Reference}
% =============================================================================

The analysis uses a panel dataset of 19 Zimbabwean self-administered pension funds
observed annually from 2018 to 2022 (95 rows total).  The file is
\texttt{new\_dataset.csv}; tests use the 10-fund subset in
\texttt{tests/sample\_data.csv}.

\begin{longtable}{>{\ttfamily}p{4.8cm} p{2.2cm} p{7.4cm}}
\toprule
\textbf{Column} & \textbf{Type} & \textbf{Description} \\
\midrule
\endhead
\bottomrule
\endfoot
fund\_id          & string  & Unique fund identifier (\texttt{ZW002}--\texttt{ZW020}).
                              Stable across years. \\[4pt]
year              & integer & Observation year: 2018, 2019, 2020, 2021, or 2022. \\[4pt]
fund\_name        & string  & Full name of the fund (e.g.\
                              \textit{Zimplats Mineworkers Pension Fund}).
                              Constant within a fund. \\[4pt]
fund\_type        & string  & \texttt{self\_administered} --- all funds in this study
                              are self-administered standalone schemes regulated by
                              IPEC. \\[4pt]
total\_assets\_usd          & float & Total assets under management expressed in
                                      USD equivalent at year-end.
                                      \textbf{DEA input.} \\[4pt]
operating\_expenses\_usd    & float & Annual administrative and management expenses
                                      in USD.  Includes staff costs, custodian fees,
                                      and audit fees.
                                      \textbf{DEA input.} \\[4pt]
equity\_debt\_usd           & float & Market value of equity plus fixed-income
                                      portfolio in USD.  Proxy for the fund's
                                      investable asset base.
                                      \textbf{DEA input.} \\[4pt]
net\_investment\_income\_usd & float & Total investment returns less investment
                                       management costs in USD.  Can be negative
                                       during crisis years (2019--2020).
                                       \textbf{DEA output.} \\[4pt]
member\_contributions\_usd  & float & Total contributions received from members and
                                      employers during the year in USD.
                                      \textbf{DEA output.} \\[4pt]
inflation         & float & Annual CPI inflation rate for Zimbabwe (\%).
                            Sourced from ZIMSTAT.
                            \textbf{RF environmental variable.} \\[4pt]
exchange\_volatility & float & Annualised standard deviation of the ZWG/USD exchange
                               rate, normalised to $[0,1]$.
                               Higher values indicate greater currency instability.
                               \textbf{RF environmental variable.} \\[4pt]
fund\_age         & integer & Years since the fund's inception as of the observation
                              year.  Increments by 1 each year.
                              \textbf{RF environmental variable.} \\
\end{longtable}

\subsection*{Derived RF features (computed in pipeline, not in CSV)}

Three additional features are derived from the panel data before the Random Forest
stage and added to the aggregated DataFrame:

\begin{table}[H]
\centering
\begin{tabular}{lll}
\toprule
\textbf{Feature} & \textbf{Formula} & \textbf{Rationale} \\
\midrule
\texttt{fund\_size\_log}  & $\ln(\text{total\_assets\_usd})$ &
    Log-transforms asset size to reduce skew. \\
\texttt{expense\_ratio}   & $\text{operating\_expenses} / \text{total\_assets}$ &
    Captures administrative cost efficiency. \\
\texttt{rts\_encoded}     & IRS$\to$1, CRS$\to$0, DRS$\to-1$ &
    Encodes returns-to-scale as an ordinal predictor. \\
\bottomrule
\end{tabular}
\caption{Derived environmental features added to the RF second stage.}
\end{table}

\subsection*{Macroeconomic variables by year}

The values below are fixed per year and identical for all funds within a year,
reflecting the shared macroeconomic environment.

\begin{table}[H]
\centering
\begin{tabular}{ccc}
\toprule
\textbf{Year} & \textbf{Inflation (\%)} & \textbf{Exchange Volatility} \\
\midrule
2018 & 10.6  & 0.18 \\
2019 & 255.3 & 0.72 \\
2020 & 557.2 & 0.85 \\
2021 & 98.5  & 0.45 \\
2022 & 243.8 & 0.61 \\
\bottomrule
\end{tabular}
\caption{Zimbabwe macroeconomic environment by year.}
\end{table}

\subsection*{Why these inputs and outputs?}

\begin{itemize}
    \item \textbf{Inputs} represent resources consumed: assets deployed, expenses
          incurred, and capital invested.  An efficient fund achieves maximum
          output from a given resource bundle.
    \item \textbf{Outputs} represent value delivered to members: investment returns
          and contribution flows.
    \item \textbf{Environmental variables} are not under the fund manager's direct
          control (inflation, exchange rate) or only slowly controllable (fund age,
          size, expense ratio, scale position).
          They are excluded from the DEA frontier but analysed in the RF second stage.
\end{itemize}

\textbf{Cross-sectional vs.\ panel:}  Before running DEA, observations are
aggregated to one row per fund by taking the mean across the five years.
This gives a single \emph{average-period} efficiency score per fund.

% =============================================================================
\section{Stage 1: CCR (CRS) DEA Model}
% =============================================================================

\subsection{Background}

Data Envelopment Analysis was introduced by Charnes, Cooper and Rhodes (1978).
The CCR model assumes Constant Returns to Scale (CRS): doubling all inputs exactly
doubles all outputs.  It measures \emph{overall technical efficiency}.

\subsection{Notation}

\begin{itemize}
    \item $n$ = number of DMUs (funds), indexed $j = 1, \ldots, n$
    \item $m$ = number of inputs, indexed $i = 1, \ldots, m$
    \item $s$ = number of outputs, indexed $r = 1, \ldots, s$
    \item $x_{ij}$ = input $i$ of DMU $j$
    \item $y_{rj}$ = output $r$ of DMU $j$
    \item $\theta_k$ = efficiency score of DMU $k$ under evaluation
    \item $\lambda_j$ = intensity variable (contribution of DMU $j$ to the reference set)
    \item $s^-_i$ = input slack for input $i$
    \item $s^+_r$ = output slack for output $r$
\end{itemize}

\subsection{Two-Phase LP Formulation}

The CCR model for DMU $k$ is solved as two sequential linear programs.

\subsubsection*{Phase I: Radial Efficiency Score}

\begin{align}
    \min_{\theta, \boldsymbol{\lambda}} \quad & \theta_k \label{eq:p1obj} \\
    \text{s.t.} \quad
    & \sum_{j=1}^{n} \lambda_j x_{ij} \leq \theta_k x_{ik}
      \quad \forall\, i = 1, \ldots, m \label{eq:p1in} \\
    & \sum_{j=1}^{n} \lambda_j y_{rj} \geq y_{rk}
      \quad \forall\, r = 1, \ldots, s \label{eq:p1out} \\
    & \lambda_j \geq 0 \quad \forall\, j, \quad 0 \leq \theta_k \leq 1 \notag
\end{align}

The constraint \eqref{eq:p1in} says the reference set must produce the evaluated
fund's outputs using no more than $\theta_k$ times its inputs.
The optimal $\theta_k^*$ is the \emph{radial contraction factor}: if $\theta_k^* = 0.8$
the fund could shrink all inputs by $20\%$ and still produce the same outputs.

\subsubsection*{Phase II: Slack Maximisation}

With $\theta_k^*$ fixed, Phase II maximises the non-radial slacks:

\begin{align}
    \max_{\boldsymbol{\lambda}, \mathbf{s}^-, \mathbf{s}^+} \quad
    & \sum_{i=1}^{m} s^-_i + \sum_{r=1}^{s} s^+_r \\
    \text{s.t.} \quad
    & \sum_{j=1}^{n} \lambda_j x_{ij} + s^-_i = \theta_k^* x_{ik}
      \quad \forall\, i \\
    & \sum_{j=1}^{n} \lambda_j y_{rj} - s^+_r = y_{rk}
      \quad \forall\, r \\
    & \lambda_j, s^-_i, s^+_r \geq 0
\end{align}

A DMU is \textbf{fully CCR-efficient} if and only if $\theta_k^* = 1$ \emph{and}
all slacks are zero.

\subsection{Why two phases?}

The single-phase formulation $\min\, \theta - \varepsilon(\sum s^-_i + \sum s^+_r)$
is theoretically equivalent but numerically unstable when input values are in the
hundreds of millions (USD).  At such scales $\varepsilon \cdot x_{ik} \gg 1$, making
the LP unbounded.  The two-phase approach avoids this entirely.

\subsection{Output file: \texttt{out/efficiency\_ccr.csv}}

\begin{table}[H]
\centering
\begin{tabular}{ll}
\toprule
\textbf{Column} & \textbf{Meaning} \\
\midrule
\texttt{fund\_id}       & Fund identifier \\
\texttt{theta\_ccr}     & CCR efficiency score $\theta_k^* \in (0,1]$ \\
\texttt{slack\_*}       & Per-input slack $s^-_i$ from Phase II \\
\texttt{peer\_ids}      & Pipe-separated IDs of reference set DMUs ($\lambda_j > 10^{-4}$) \\
\bottomrule
\end{tabular}
\end{table}

% =============================================================================
\section{Stage 2: BCC (VRS) DEA Model}
% =============================================================================

\subsection{Background}

The BCC model (Banker, Charnes, Cooper, 1984) adds a \emph{convexity constraint}
$\sum_j \lambda_j = 1$, relaxing the CRS assumption to allow Variable Returns to
Scale.  BCC measures \emph{pure technical efficiency}, isolating managerial
efficiency from scale effects.

\subsection{Modification to Phase I}

The only change to the CCR Phase I LP is the addition of:
\[
    \sum_{j=1}^{n} \lambda_j = 1
\]

This forces the reference virtual DMU to lie on the convex hull of the observed
data rather than on a ray through the origin.

\begin{remark}
    Because VRS is more permissive than CRS, BCC scores are always
    $\geq$ CCR scores for the same DMU: $\theta_k^{\text{BCC}} \geq \theta_k^{\text{CCR}}$.
\end{remark}

\subsection{Output file: \texttt{out/efficiency\_vrs.csv}}

Same structure as \texttt{efficiency\_ccr.csv} but the score column is
\texttt{theta\_bcc}.

% =============================================================================
\section{Stage 3: Scale Efficiency and Returns to Scale}
% =============================================================================

\subsection{Scale efficiency}

\begin{definition}
    The \textbf{scale efficiency} of DMU $k$ is:
    \[
        SE_k = \frac{\theta_k^{\text{CCR}}}{\theta_k^{\text{BCC}}}
    \]
\end{definition}

$SE_k \in (0, 1]$.  A value of 1 means the fund is operating at its most
productive scale size (MPSS); values below 1 indicate scale inefficiency.

\subsection{Returns to Scale classification}

The nature of returns to scale is determined from the BCC optimal lambda vector
$\boldsymbol{\lambda}^*$:

\begin{table}[H]
\centering
\begin{tabular}{lll}
\toprule
\textbf{Condition} & \textbf{Classification} & \textbf{Interpretation} \\
\midrule
$\sum_j \lambda_j^* \approx 1$ and $SE_k \approx 1$ & CRS & Optimal scale \\
$\sum_j \lambda_j^* < 1$                              & IRS & Fund is too small \\
$\sum_j \lambda_j^* > 1$                              & DRS & Fund is too large \\
\bottomrule
\end{tabular}
\end{table}

\textbf{Expected pattern in our data:}
\begin{itemize}
    \item Large long-established funds (e.g.\ Delta Corporation, POSB) $\rightarrow$ DRS
    \item Small younger funds (e.g.\ Cottco, NetOne) $\rightarrow$ IRS
    \item Medium-sized funds operating close to their optimal scale $\rightarrow$ CRS
\end{itemize}

The RTS classification is also encoded numerically as \texttt{rts\_encoded}
(IRS$\to$1, CRS$\to$0, DRS$\to-1$) and used as a predictor in the RF second stage.

\subsection{Output file: \texttt{out/scale.csv}}

Columns: \texttt{fund\_id}, \texttt{scale\_efficiency}, \texttt{rts\_classification}.

% =============================================================================
\section{Stage 4: Peer Analysis and Input Reduction Targets}
% =============================================================================

\subsection{Target inputs}

For each inefficient DMU $k$, the \emph{projected} (target) inputs are:
\[
    x_{ik}^* = \theta_k^* \cdot x_{ik} - s^{-*}_{ik}
\]
where $s^{-*}_{ik}$ is the Phase II slack for input $i$.

\subsection{Reduction metrics}

\[
    \Delta x_{ik} = x_{ik} - x_{ik}^*
    \qquad
    \Delta x_{ik}^{\%} = \frac{\Delta x_{ik}}{x_{ik}} \times 100
\]

These tell the fund manager: ``to reach the frontier, you must reduce
\textit{operating expenses} by \$X (Y\%) and \textit{total assets deployed} by \$A (B\%).''

\subsection{Output file: \texttt{out/targets.csv}}

Columns: \texttt{fund\_id}, \texttt{theta\_ccr}, then for each input:
\texttt{target\_*}, \texttt{reduction\_abs\_*}, \texttt{reduction\_pct\_*}.

% =============================================================================
\section{Stage 5: Simar--Wilson Bootstrap Bias Correction}
% =============================================================================

\subsection{Why bias correction is needed}

The DEA estimator $\hat{\theta}_k$ is upward biased: the finite-sample frontier
lies inside the true population frontier, inflating efficiency scores.
The Simar--Wilson (1998) Algorithm 1 corrects this bias using a smoothed bootstrap.

\subsection{Algorithm}

Let $\hat{\boldsymbol{\theta}} = (\hat{\theta}_1, \ldots, \hat{\theta}_n)$ be the
original DEA scores.

\textbf{Step 1 -- Silverman bandwidth.}
\[
    h = 0.9 \cdot \min\!\left(\hat{\sigma},\, \frac{\text{IQR}(\hat{\boldsymbol{\theta}})}{1.34}\right)
    \cdot n^{-1/5}
\]

\textbf{Step 2 -- Reflected KDE resampling.}  For $b = 1, \ldots, B$:
\begin{enumerate}[label=\alph*.]
    \item Draw $\hat{\theta}_1^*, \ldots, \hat{\theta}_n^*$ i.i.d.\ from the kernel
          density of $\hat{\boldsymbol{\theta}}$ with bandwidth $h$.
    \item Reflect values above 1 back: $\hat{\theta}_j^* \leftarrow 2 - \hat{\theta}_j^*$
          if $\hat{\theta}_j^* > 1$.
    \item Clip to $(0, 1]$.
\end{enumerate}

\textbf{Step 3 -- Pseudo-dataset.}  Rescale inputs to impose pseudo-efficiency:
\[
    x_{ij}^{*(b)} = \frac{\hat{\theta}_j^{*(b)}}{\hat{\theta}_j} \cdot x_{ij}
\]

\textbf{Step 4 -- Bootstrap DEA.}  Solve DEA on $(X^{*(b)}, Y)$ to get
$\hat{\theta}_k^{*(b)}$.

\textbf{Step 5 -- Bias and corrected score.}
\[
    \widehat{\text{bias}}_k = \frac{1}{B}\sum_{b=1}^{B} \hat{\theta}_k^{*(b)} - \hat{\theta}_k
    \qquad
    \tilde{\theta}_k = \hat{\theta}_k - \widehat{\text{bias}}_k
\]

\textbf{Step 6 -- Confidence interval.}  The $95\%$ CI is the $[2.5\%, 97.5\%]$
percentiles of $\{\hat{\theta}_k^{*(b)}\}_{b=1}^{B}$.

The default is $B = 2000$; $B = 200$ is used in CI/test runs.
Bootstrap replications are parallelised with \texttt{joblib}.

\subsection{Output file: \texttt{out/bias\_corrected\_scores.csv}}

Columns: \texttt{fund\_id}, \texttt{theta\_raw}, \texttt{theta\_bias\_corrected},
\texttt{bias}, \texttt{ci\_lower}, \texttt{ci\_upper}.

% =============================================================================
\section{Stage 6: Random Forest Second Stage}
% =============================================================================

\subsection{Purpose}

The second-stage regression identifies which environmental and structural factors
explain variation in bias-corrected efficiency scores $\tilde{\theta}_k$.
Six candidate determinants are used, combining macroeconomic variables, operational
ratios, and the RTS classification derived from the DEA itself.

\subsection{Model}

\[
    \tilde{\theta}_k = f\!\bigl(
        \text{inflation}_k,\;
        \text{exchange\_volatility}_k,\;
        \text{fund\_age}_k,\;
        \ln(\text{assets}_k),\;
        \text{expense\_ratio}_k,\;
        \text{rts\_encoded}_k
    \bigr) + \varepsilon_k
\]

$f(\cdot)$ is estimated by a Random Forest with:
\begin{itemize}
    \item $n\_\text{estimators} = 500$ trees
    \item $\text{max\_features} = \sqrt{p}$ features per split (Breiman's default)
    \item Deterministic via \texttt{np.random.default\_rng(seed)}
    \item 5-fold cross-validation to assess out-of-sample $R^2$
\end{itemize}

\subsection{Derived feature construction (in \texttt{cli.py}, Stage 7 block)}

Before fitting the RF model, the aggregated DataFrame is augmented with:
\begin{align*}
    \texttt{fund\_size\_log}_k &= \ln(\overline{\text{total\_assets}}_{k}) \\
    \texttt{expense\_ratio}_k &= \frac{\overline{\text{operating\_expenses}}_{k}}
                                      {\overline{\text{total\_assets}}_{k}} \\
    \texttt{rts\_encoded}_k &=
        \begin{cases} +1 & \text{if IRS} \\ 0 & \text{if CRS} \\ -1 & \text{if DRS} \end{cases}
\end{align*}
where bars denote the five-year mean.

\subsection{Feature importance measures}

\textbf{MDI importance (impurity-based):}
\[
    I_j^{\text{MDI}} = \frac{1}{T} \sum_{t=1}^{T} \sum_{\text{node } v \in t:\, \text{split on } j}
    \frac{N_v}{N} \Delta \text{impurity}(v)
\]

\textbf{Permutation importance:}  Feature $j$ is permuted 30 times; the mean
decrease in $R^2$ is recorded.  This avoids MDI's bias toward high-cardinality
features.

\subsection{Partial Dependence Plots}

The marginal effect of feature $j$ on the predicted score, averaging over all
other features:
\[
    \hat{f}_j(x_j) = \frac{1}{n} \sum_{i=1}^{n} \hat{f}(x_j, \mathbf{x}_{i, -j})
\]

Saved as \texttt{out/pdp\_top3.png} for the three most important features.

\subsection{Output files}

\begin{itemize}
    \item \texttt{out/rf\_importance.csv} --- feature name, MDI importance
    \item \texttt{out/pdp\_top3.png}     --- partial dependence plots
\end{itemize}

% =============================================================================
\section{Graph Descriptions}
% =============================================================================

\subsection{CCR vs.\ BCC Efficiency Bar Chart}

Side-by-side bars per fund.  BCC bars are always $\geq$ CCR bars.
\textbf{What to look for:}
\begin{itemize}
    \item Funds where CCR $\approx$ BCC: scale-efficient.
    \item Large gap (BCC $\gg$ CCR): scale inefficiency is the dominant problem.
    \item Low BCC: pure managerial inefficiency regardless of scale.
\end{itemize}

\subsection{Scale Efficiency Chart}

Bar chart of $SE_k$ per fund, coloured by RTS class (IRS / CRS / DRS).
Funds below $SE = 0.9$ are flagged for scale restructuring.

\subsection{Partial Dependence Plots}

Three sub-plots, one per top driver.  The y-axis is the predicted
$\tilde{\theta}$ and the x-axis is the feature value.  A downward slope for
\textit{inflation} indicates that higher inflation reduces efficiency;
an upward slope for \textit{fund\_age} suggests older funds are more efficient
(institutional learning).  The \textit{expense\_ratio} plot shows whether
cost-heavy funds systematically underperform.

% =============================================================================
\section{Process Verification: Python Test Snippets}
% =============================================================================

The following snippets test the correctness of each pipeline stage.
Run them with \texttt{uv run pytest explanation/test\_process.py -v}
after placing \texttt{test\_process.py} alongside this file,
or run interactively by pasting into a REPL inside the project directory.

\bigskip

\subsection{Stage 0: Data loading and schema}

\begin{lstlisting}[style=python, caption={Test: CSV loads with correct schema and types}]
import pandas as pd
from pathlib import Path

def test_dataset_schema():
    df = pd.read_csv("tests/sample_data.csv")

    # Correct number of rows and unique funds
    assert len(df) == 20, f"Expected 20 rows, got {len(df)}"
    assert df["fund_id"].nunique() == 10, "Expected 10 unique funds"
    assert sorted(df["year"].unique()) == [2021, 2022]

    # No (fund_id, year) duplicates
    assert df.duplicated(["fund_id", "year"]).sum() == 0

    # All funds are self-administered
    assert (df["fund_type"] == "self_administered").all(), \
        "All funds must be self_administered"

    # Financial columns must be positive (except net_investment_income_usd)
    for col in ["total_assets_usd", "operating_expenses_usd",
                "equity_debt_usd", "member_contributions_usd"]:
        assert (df[col] > 0).all(), f"{col} has non-positive values"

    print("PASS: dataset schema OK")

test_dataset_schema()
\end{lstlisting}

\subsection{Stage 1: CCR DEA correctness}

\begin{lstlisting}[style=python, caption={Test: CCR scores are in (0,1] and frontier funds score 1}]
import numpy as np
import pandas as pd
from pension_toolkit.data_io import load_csv, get_dea_matrices
from pension_toolkit.dea_core import dea_ccr_input_oriented

def test_ccr_dea():
    df = load_csv("tests/sample_data.csv")
    numeric_cols = df.select_dtypes("number").columns.tolist()
    df_agg = df.groupby("fund_id")[numeric_cols].mean().reset_index()

    INPUT_COLS  = ["operating_expenses_usd", "total_assets_usd", "equity_debt_usd"]
    OUTPUT_COLS = ["net_investment_income_usd", "member_contributions_usd"]
    fund_ids = df_agg["fund_id"].tolist()
    X = df_agg[INPUT_COLS].to_numpy(dtype=float)
    Y = df_agg[OUTPUT_COLS].to_numpy(dtype=float)

    result = dea_ccr_input_oriented(X, Y, fund_ids)

    # All scores must be in (0, 1]
    assert (result.theta > 0).all(),  "Zero efficiency score detected"
    assert (result.theta <= 1 + 1e-6).all(), "Score exceeds 1.0"

    # At least one fully efficient fund (theta == 1)
    assert (result.theta >= 1 - 1e-4).any(), \
        "No fund is on the CCR frontier -- possible data problem"

    # Not all funds are efficient (otherwise DEA is trivial)
    assert (result.theta < 0.999).any(), "All funds score 1.0 -- no inefficiency detected"

    # Lambda matrix shape
    n = len(fund_ids)
    assert result.lambdas.shape == (n, n)
    assert result.slacks_in.shape == (n, len(INPUT_COLS))

    print(f"PASS: CCR mean={result.theta.mean():.4f}, "
          f"efficient={( result.theta >= 1-1e-4).sum()}/{n}")

test_ccr_dea()
\end{lstlisting}

\subsection{Stage 2: BCC $\geq$ CCR and Scale Efficiency in (0, 1]}

\begin{lstlisting}[style=python, caption={Test: BCC dominance and scale efficiency bounds}]
import numpy as np
import pandas as pd
from pension_toolkit.data_io import load_csv, get_dea_matrices
from pension_toolkit.dea_core import dea_ccr_input_oriented, dea_bcc_input_oriented
from pension_toolkit.scale import compute_scale_efficiency

def test_scale_efficiency():
    df = load_csv("tests/sample_data.csv")
    numeric_cols = df.select_dtypes("number").columns.tolist()
    df_agg = df.groupby("fund_id")[numeric_cols].mean().reset_index()

    INPUT_COLS  = ["operating_expenses_usd", "total_assets_usd", "equity_debt_usd"]
    OUTPUT_COLS = ["net_investment_income_usd", "member_contributions_usd"]
    fund_ids = df_agg["fund_id"].tolist()
    X = df_agg[INPUT_COLS].to_numpy(dtype=float)
    Y = df_agg[OUTPUT_COLS].to_numpy(dtype=float)

    ccr = dea_ccr_input_oriented(X, Y, fund_ids)
    bcc = dea_bcc_input_oriented(X, Y, fund_ids)

    # BCC >= CCR for all DMUs
    assert (bcc.theta >= ccr.theta - 1e-6).all(), "BCC < CCR for some DMU"

    scale = compute_scale_efficiency(ccr, bcc)

    # Scale efficiency in (0, 1]
    assert (scale.scale_efficiency > 0).all()
    assert (scale.scale_efficiency <= 1 + 1e-6).all()

    # RTS classification is one of IRS / DRS / CRS
    valid_rts = {"IRS", "DRS", "CRS"}
    assert all(r in valid_rts for r in scale.rts_class)

    print(f"PASS: mean SE={scale.scale_efficiency.mean():.4f}, "
          f"IRS={scale.rts_class.count('IRS')}, "
          f"DRS={scale.rts_class.count('DRS')}, "
          f"CRS={scale.rts_class.count('CRS')}")

test_scale_efficiency()
\end{lstlisting}

\subsection{Stage 3: Input reduction targets are valid}

\begin{lstlisting}[style=python, caption={Test: target inputs are non-negative and <= actual inputs}]
import pandas as pd

def test_targets():
    targets = pd.read_csv("out/targets.csv")
    ccr     = pd.read_csv("out/efficiency_ccr.csv")
    dataset = pd.read_csv("tests/sample_data.csv")

    numeric = dataset.select_dtypes("number").columns.tolist()
    df_agg  = dataset.groupby("fund_id")[numeric].mean().reset_index()

    INPUT_COLS = ["operating_expenses_usd", "total_assets_usd", "equity_debt_usd"]

    for col in INPUT_COLS:
        target_col = f"target_{col}"
        if target_col not in targets.columns:
            continue
        # Targets must be non-negative
        assert (targets[target_col] >= 0).all(), f"Negative target for {col}"
        # Targets must be <= actual
        actual = df_agg.set_index("fund_id")[col]
        for _, row in targets.iterrows():
            act = actual.loc[row["fund_id"]]
            assert row[target_col] <= act + 1e-3, \
                f"Target > actual for {row['fund_id']} / {col}"

    # Efficient funds (theta==1) should have zero reduction
    efficient = ccr[ccr["theta_ccr"] >= 1 - 1e-4]["fund_id"].tolist()
    for fid in efficient:
        row = targets[targets["fund_id"] == fid]
        for col in INPUT_COLS:
            pct_col = f"reduction_pct_{col}"
            if pct_col in row.columns:
                assert row[pct_col].values[0] < 1.0, \
                    f"Efficient fund {fid} has non-zero reduction for {col}"

    print(f"PASS: targets valid for all {len(targets)} funds")

test_targets()
\end{lstlisting}

\subsection{Stage 4: Bootstrap bias correction properties}

\begin{lstlisting}[style=python, caption={Test: bias-corrected scores are in [0,1]; CIs are valid}]
import numpy as np
import pandas as pd
from pension_toolkit.data_io import load_csv, get_dea_matrices
from pension_toolkit.dea_core import dea_ccr_input_oriented
from pension_toolkit.bootstrap import simar_wilson

def test_bootstrap():
    df = load_csv("tests/sample_data.csv")
    numeric = df.select_dtypes("number").columns.tolist()
    df_agg = df.groupby("fund_id")[numeric].mean().reset_index()

    INPUT_COLS  = ["operating_expenses_usd", "total_assets_usd", "equity_debt_usd"]
    OUTPUT_COLS = ["net_investment_income_usd", "member_contributions_usd"]
    fund_ids = df_agg["fund_id"].tolist()
    X = df_agg[INPUT_COLS].to_numpy(dtype=float)
    Y = df_agg[OUTPUT_COLS].to_numpy(dtype=float)

    def ccr_func(X_, Y_):
        return dea_ccr_input_oriented(X_, Y_, fund_ids)

    result = simar_wilson(ccr_func, X, Y, fund_ids=fund_ids, B=200, seed=42)

    # Bias-corrected scores are in [0, 1]
    assert (result.theta_bias_corrected >= 0).all()
    assert (result.theta_bias_corrected <= 1 + 1e-6).all()

    # CIs are ordered: lower <= upper
    assert (result.ci_lower <= result.ci_upper + 1e-9).all()

    # Reproducibility
    r2 = simar_wilson(ccr_func, X, Y, fund_ids=fund_ids, B=200, seed=42)
    np.testing.assert_array_almost_equal(
        result.theta_bias_corrected, r2.theta_bias_corrected,
        decimal=6, err_msg="Bootstrap is not reproducible with same seed"
    )

    print(f"PASS: mean raw={result.theta_raw.mean():.4f}, "
          f"mean BC={result.theta_bias_corrected.mean():.4f}, "
          f"mean bias={result.bias.mean():.4f}")

test_bootstrap()
\end{lstlisting}

\subsection{Stage 5: Derived features and Random Forest second stage}

\begin{lstlisting}[style=python, caption={Test: derived features correct; RF importance sums to 1}]
import numpy as np
import pandas as pd
from pension_toolkit.data_io import load_csv
from pension_toolkit.dea_core import dea_ccr_input_oriented, dea_bcc_input_oriented
from pension_toolkit.scale import compute_scale_efficiency
from pension_toolkit.ml_stage import fit_rf, ENV_FEATURE_COLS

def test_rf_second_stage():
    df = load_csv("tests/sample_data.csv")
    numeric = df.select_dtypes("number").columns.tolist()
    df_agg = df.groupby("fund_id")[numeric].mean().reset_index()

    INPUT_COLS  = ["operating_expenses_usd", "total_assets_usd", "equity_debt_usd"]
    OUTPUT_COLS = ["net_investment_income_usd", "member_contributions_usd"]
    fund_ids = df_agg["fund_id"].tolist()
    X = df_agg[INPUT_COLS].to_numpy(dtype=float)
    Y = df_agg[OUTPUT_COLS].to_numpy(dtype=float)

    # Build bias-corrected scores inline (small B for speed)
    from pension_toolkit.bootstrap import simar_wilson
    def ccr_func(X_, Y_): return dea_ccr_input_oriented(X_, Y_, fund_ids)
    boot = simar_wilson(ccr_func, X, Y, fund_ids=fund_ids, B=50, seed=42)

    # Compute derived RF features (mirrors cli.py logic)
    ccr = dea_ccr_input_oriented(X, Y, fund_ids)
    bcc = dea_bcc_input_oriented(X, Y, fund_ids)
    scale = compute_scale_efficiency(ccr, bcc)
    rts_map = {"IRS": 1, "CRS": 0, "DRS": -1}
    df_agg["fund_size_log"] = np.log(df_agg["total_assets_usd"])
    df_agg["expense_ratio"] = (df_agg["operating_expenses_usd"]
                                / df_agg["total_assets_usd"])
    df_agg["rts_encoded"] = [rts_map.get(r, 0) for r in scale.rts_class]

    env_cols = [c for c in ENV_FEATURE_COLS if c in df_agg.columns]
    env_features = df_agg[env_cols].values.astype(float)

    result = fit_rf(boot.theta_bias_corrected, env_features,
                    feature_names=env_cols, seed=42)

    # MDI importance sums to 1
    total_imp = result.feature_importance["importance"].sum()
    assert abs(total_imp - 1.0) < 1e-6, f"Importance sum = {total_imp}"

    # All feature names accounted for
    assert set(result.feature_importance["feature"]) == set(env_cols)

    # top3 contains valid feature names
    assert all(f in env_cols for f in result.top3_features)

    print(f"PASS: CV R2 mean={result.cv_r2_scores.mean():.4f}, "
          f"top driver={result.top3_features[0]}")

test_rf_second_stage()
\end{lstlisting}

\subsection{Stage 6: End-to-end pipeline output files}

\begin{lstlisting}[style=python, caption={Test: all output files are produced and non-empty}]
from pathlib import Path
import pandas as pd

EXPECTED_OUTPUTS = [
    ("out/efficiency_ccr.csv",        ["fund_id", "theta_ccr"]),
    ("out/efficiency_vrs.csv",        ["fund_id", "theta_bcc"]),
    ("out/scale.csv",                 ["fund_id", "scale_efficiency",
                                       "rts_classification"]),
    ("out/targets.csv",               ["fund_id", "theta_ccr"]),
    ("out/bias_corrected_scores.csv", ["fund_id", "theta_raw",
                                       "theta_bias_corrected",
                                       "ci_lower", "ci_upper"]),
    ("out/rf_importance.csv",         ["feature", "importance"]),
]

def test_pipeline_outputs():
    for csv_path, required_cols in EXPECTED_OUTPUTS:
        p = Path(csv_path)
        assert p.exists(), f"Missing output: {csv_path}"
        df = pd.read_csv(p)
        assert len(df) > 0, f"Empty file: {csv_path}"
        for col in required_cols:
            assert col in df.columns, f"Column '{col}' missing from {csv_path}"

    assert Path("out/report.pdf").exists(), "PDF report not generated"
    assert Path("out/pdp_top3.png").exists(), "PDP plot not generated"

    print("PASS: all output files present with required columns")

test_pipeline_outputs()
\end{lstlisting}

% =============================================================================
\section{Running the Full Pipeline}
% =============================================================================

\begin{lstlisting}[style=python, language=bash, caption={Shell commands to reproduce all results}]
# 1. Enter project directory and sync environment
cd prototype/pension_efficiency_toolkit
uv sync

# 2. Run full analysis on the real dataset (200 bootstrap replications for speed)
uv run python -m pension_toolkit.cli analyze \
    --input ../../new_dataset.csv \
    --out out/ \
    --bootstrap-B 200 \
    --seed 42

# 3. Run unit tests
uv run pytest -q

# 4. Run process-verification snippets
uv run python explanation/test_process.py

# 5. Launch dashboard
uv run streamlit run pension_toolkit/ui_streamlit.py
\end{lstlisting}

% =============================================================================
\section{Key References}
% =============================================================================

\begin{enumerate}
    \item Charnes, A., Cooper, W.W., Rhodes, E.\ (1978). Measuring the efficiency
          of decision making units. \textit{European Journal of Operational Research},
          2(6), 429--444.
    \item Banker, R.D., Charnes, A., Cooper, W.W.\ (1984). Some models for estimating
          technical and scale inefficiencies in data envelopment analysis.
          \textit{Management Science}, 30(9), 1078--1092.
    \item Simar, L., Wilson, P.W.\ (1998). Sensitivity analysis of efficiency scores:
          How to bootstrap in nonparametric frontier models.
          \textit{Management Science}, 44(1), 49--61.
    \item Breiman, L.\ (2001). Random forests. \textit{Machine Learning}, 45(1), 5--32.
    \item Insurance and Pensions Commission Zimbabwe (IPEC). Annual Reports 2018--2022.
\end{enumerate}

\end{document}
